% ============================================================================
% 问题一 —— 三段式：分析与准备 -> 建模与求解 -> 结果检验。
%
% 【四问通用的写作合同】
%   1. 先出数再写文：本文件里每个数字都来自 numbers.tex 宏，写作时台账必须
%      已经冻结。禁止先写空论文再填数。
%   2. 主结果旁必须有证书：置信下界 / 收敛判据 / 双路互证三者之一，
%      只给点估计的结论一律降措辞。
%   3. 最优性主张分级：global / local / best-found 三选一，
%      无据不得写「最优」「最低」「证明」。
%   4. 每问至少一张由代码生成的关键图，图与正文数字必须对得上。
%   5. 推导过程移附录，正文只留「设定 -> 关键结论 -> 检验」这条证据链主线。
%
% 写长了（超过 6 页）再拆片：05a-问题一-分析与准备.tex /
% 05b-问题一-建模与求解.tex，main.tex 里把一行 \input 换成两行即可——
% 分节粒度本来就是可调的。
% ============================================================================
\section{问题一模型的建立与求解}

\subsection{分析与准备}

\subsection{模型的建立}

\subsection{模型的求解}

\subsection{结果与检验}
